% ============================================================================
% Fichier  : front/02_mode_emploi.tex
% Titre    : Mode d'emploi du cours
% Auteur   : MINKA MI NGUIDJOÏ Thierry Emmanuel
% Version  : 1.0 -- novembre 2026
% ============================================================================

\chapter*{Mode d'emploi du cours}
\addcontentsline{toc}{chapter}{Mode d'emploi du cours}
\markboth{Mode d'emploi}{Mode d'emploi}

\section*{Structure générale}

Le cours compte un prologue et six parties, soit 26~chapitres, pour
60~heures en présentiel réparties en trente séances de deux heures. Les
fiches d'auto-évaluation et les dix laboratoires de NDAP SUKULU sont des
travaux à faire hors de ces 60~heures.

\begin{tableaucours}{L{1.6cm}L{4.4cm}L{4.4cm}L{1.8cm}}{%
    \tableauHeader{Partie} &
    \tableauHeader{Thème} &
    \tableauHeader{Compétences visées} &
    \tableauHeader{Niveau}
}
Prologue & Une nuit à Garoua & Lire un incident, séparer constatation et
opinion & \niveaubase \\
I & Fondements & Propriétés, menaces, cryptographie, histoire &
\niveaubase \\
II & Gouvernance et risques & Rôles, politique, risques, actifs, facteur
humain & \niveaubase \\
III & Normes et droit & Référentiels, loi camerounaise, éthique &
\niveaubase \\
IV & Sécurité technique & Architecture, réseaux, identités, systèmes,
applications, nuage, opérations & \niveaumoyen \\
V & Défis avancés & Continuité, investigation, technologies émergentes &
\niveaumoyen{} à \niveauexpert \\
VI & Mesurer & Indicateurs, \OUI{}~/ \NON{}~/ \NEK{} & \niveauexpert \\
\end{tableaucours}

\section*{Les trente séances}

\begin{tableaucours}{L{1.4cm}L{8.4cm}L{2.4cm}}{%
    \tableauHeader{Séance} &
    \tableauHeader{Contenu} &
    \tableauHeader{Partie}
}
1 & Prologue~: une nuit à Garoua & Prologue \\
2 & Objets et propriétés (Ch.~\ref{chap:proprietes}) & I \\
3 & Menaces et attaques (Ch.~\ref{chap:menaces}) & I \\
4 & Cryptographie symétrique (Ch.~\ref{chap:crypto-sym}) & I \\
5 & Cryptographie asymétrique (Ch.~\ref{chap:crypto-asym}) & I \\
6 & Grandes attaques (Ch.~\ref{chap:attaques}) & I \\
7 & Gouvernance (Ch.~\ref{chap:gouvernance}) & II \\
8 & Risques (Ch.~\ref{chap:risques}) & II \\
9 & Actifs et données (Ch.~\ref{chap:actifs}) & II \\
10 & Facteur humain (Ch.~\ref{chap:humain}) & II \\
11 & Normes et référentiels (Ch.~\ref{chap:normes}) & III \\
12 & Droit (Ch.~\ref{chap:droit}) & III \\
13 & Éthique (Ch.~\ref{chap:ethique}) & III \\
14 & Modèles formels (Ch.~\ref{chap:modeles}) & IV \\
15 & Architecture (Ch.~\ref{chap:architecture}) & IV \\
16 & Réseaux~: protocoles et attaques (Ch.~\ref{chap:reseaux1}) & IV \\
17 & Réseaux~: équipements (Ch.~\ref{chap:reseaux2}) & IV \\
18--19 & Identités et accès (Ch.~\ref{chap:iam}) & IV \\
20 & Sécurité des systèmes (Ch.~\ref{chap:systemes}) & IV \\
21--22 & Sécurité des applications (Ch.~\ref{chap:applications}) & IV \\
23 & Nuage et mobiles (Ch.~\ref{chap:nuage}) & IV \\
24 & Évaluer la sécurité (Ch.~\ref{chap:evaluer}) & IV \\
25--26 & Opérations et incidents (Ch.~\ref{chap:operations}) & IV \\
27 & Continuité et reprise (Ch.~\ref{chap:continuite}) & V \\
28 & Investigation numérique (Ch.~\ref{chap:investigation}) & V \\
29 & Défis émergents (Ch.~\ref{chap:emergents}) & V \\
30 & Mesurer la sécurité~; synthèse (Ch.~\ref{chap:mesurer}) & VI \\
\end{tableaucours}

\section*{Les encadrés}

\begin{boxdef}
\textbf{Définition}. Introduit un terme avec sa définition~; les termes
définis sont indexés et repris dans le glossaire (Annexe~\ref{ann:glossaire}).
\end{boxdef}

\begin{boxscenario}{Titre du scénario}
\textbf{Scénario opérationnel}. Un épisode d'une des opérations du
cours~: GAROUA et LIZBIZ (LiZbiZ), WOURI (BANQUE SANAGA), MEFOU (ANPP).
\end{boxscenario}

\begin{boxobjet}{M-xx · Intitulé}
\textbf{Mesure de sécurité}. Présente une mesure sous une forme
constante~: objectif, menace couverte, mise en œuvre, vérification.
\end{boxobjet}

\begin{boxprocedure}
\textbf{Procédure opérationnelle}. Dit, pas à pas, ce qu'il faut faire.
\end{boxprocedure}

\begin{boxmetier}{RSSI}
\textbf{Application métier}. Traduit le chapitre en gestes concrets pour
l'un des métiers visés~: RSSI, administrateur systèmes et réseaux, analyste
SOC, développeur, auditeur de sécurité.
\end{boxmetier}

\begin{boxjuris}
Cite un texte de loi d'après sa version officielle, puis en tire la
portée pour la sécurité.
\end{boxjuris}

\begin{boxalert}
\textbf{Violation typique} et \textbf{règle}~: une erreur fréquente, et ce
qu'il faut faire à la place.
\end{boxalert}

\begin{boxexo}
\textbf{Exercice}. Dix par chapitre, gradués en trois niveaux~:
\niveaubase\ fondamental, \niveaumoyen\ intermédiaire, \niveauexpert\
expert. Les corrigés figurent dans le livret de correction.
\end{boxexo}

\section*{Conventions}

\begin{itemize}
    \item Chaque analyse indique un niveau de confiance~: faible, moyenne ou
          élevée, et reste séparée de la constatation.
    \item Les \termecle{termes clés} apparaissent en gras~; les normes
          s'écrivent \norme{ISO/IEC~27001}, les lois camerounaises
          \loicam{n°~2010/012}.
    \item Les seuils et les échelles proposés (cotation des risques, délais
          de correctifs) sont des conventions pédagogiques, à adapter par
          chaque organisation.
\end{itemize}

\section*{Matériel d'accompagnement}

\begin{itemize}
    \item \textbf{Annexes}~: glossaire, outils, modèles et procédures,
          textes juridiques et normatifs, bibliographie commentée, contacts
          (Annexes~\ref{ann:glossaire} à~\ref{ann:contacts}).
    \item \textbf{Livret de correction}~: corrigés des exercices et des
          fiches, réservé au formateur.
    \item \textbf{NDAP SUKULU}~: fiches d'auto-évaluation et laboratoires,
          présentés dans la section qui suit.
\end{itemize}
